\chapter{INTRODUÇÃO}
\label{capitulo-introducao}

Este projeto de extensão propõe o desenvolvimento de um protótipo de bracelete inteligente utilizando a plataforma ESP32, com a finalidade de auxiliar na orientação e na percepção por pessoas com baixa visão. A proposta relaciona conhecimentos de programação, eletrônica, sistemas embarcados e Internet das Coisas (IoT) com uma necessidade de caráter social, buscando aplicar a computação no desenvolvimento de uma tecnologia assistiva de baixo custo.

O público-alvo da proposta são pessoas com baixa visão que possam se beneficiar de um dispositivo vestível de auxílio à percepção do ambiente. A atividade será desenvolvida no Instituto Federal do Paraná – Campus Pinhais, em ambiente acadêmico apropriado para montagem, programação e testes do protótipo. Os extensionistas serão os estudantes responsáveis pelo projeto, com acompanhamento e orientação do professor responsável pela disciplina.

